\beginsong{Et toi compagnon !}[by={Bertrand / Louis et Elisa — SGDF Rhône Ouest}]
\beginchorus
Et t\[Bm]oi compagnon, que mets-tu dans ton \[A]sac ?
Quitte ta maison, quitte ton ham\[Bm]ac 
Pour quelles destinat\[A]ions, vers quel tarm\[F#]ac
Et t\[Bm]oi compagnon, que mets-tu dans ton \[A]sac ?
Quitte ta maison, quitte ton ham\[Bm]ac 
Pour quelles acti\[A]ons, vers quel bivou\[F#]ac
\endchorus
\beginverse
\[G]Découvrir au fond de \[D]toi, la beauté de tes rencontr\[C]es,
Ne pas les garder pour \[D]toi, il faut que tu nous les montr\[G]es,
À l'écoute de ton c{\oe}\[D]ur, à l'affût d'un beau souri\[C]re,
Tu seras toujours à l'h\[D]eure, tu sauras nous réu\[F#]nir !
\endverse
\beginverse
\[G]Découvrir au fond de \[D]toi, ce que tu veux deve\[C]nir,
Toujours tu auras la \[D]foi, sur la route, ton ave\[G]nir,
Ranger ces richesses ferti\[D]les, il faudra les faire pous\[C]ser,
Tu n'seras pas malhab\[D]ile, il faudra les parta\[F#]ger !
\endverse
\beginverse
\[G]Découvrir au fond de \[D]toi, ces paroles de l'amit\[C]ié,
Savoir être proche, être \[D]là, pour des liens d'éterni\[G]té,
Rester souvent à l'éco\[D]ute, tel une étoile, un berg\[C]er,
Compa de routes et de dout\[D]es, tu sauras comment parl\[F#]er !
\endverse
\beginverse
\[G]Découvrir au fond de \[D]toi, tous ces acteurs solidair\[C]es,
Citoyen en qui l'on cr\[D]oit, tu les mettras en lumiè\[G]re,
Grandir en témoin du mon\[D]de, agir en réalit\[C]é,
Toujours prêt comme il l'a \[D]dit, pour cette terre, oui t'enga\[F#]ger !
\endverse
\endsong
